
\documentclass[11pt]{article}

% ============================================================
% Working title:
% GeoODE-KD: Distilling Sentence Embeddings as
% Teacher-Guided Geometric Dynamics
%
% Suggested packages:
\usepackage{amsmath,amssymb,amsfonts}
\usepackage{booktabs}
\usepackage{algorithm}
\usepackage{algpseudocode}
\usepackage[round]{natbib}
\usepackage[margin=1in]{geometry}
\usepackage{microtype}
% ============================================================


\title{GeoODE-KD\@: Distilling Sentence Embeddings as\\
Teacher-Guided Geometric Dynamics}
\author{}
\date{}

\begin{document}

\maketitle


% ============================================================
% ABSTRACT
% ============================================================

\begin{abstract}
Knowledge distillation provides an effective route for transferring the
semantic capabilities of large text embedding models into compact encoders.
Existing embedding distillation methods predominantly formulate knowledge
transfer as \emph{state matching}: the student is encouraged to reproduce
teacher embeddings, intermediate representations, or pairwise relations at
one or more discrete network depths. Such objectives specify where student
representations should be located, but provide little structure on
\emph{how} they should evolve through depth. This limitation becomes
particularly relevant under large teacher-student capacity gaps, where
forcing shallow student layers toward highly abstract teacher
representations can over-constrain their native feature organization.

We propose \textbf{GeoODE-KD}, a continuous-depth formulation of sentence
embedding distillation that transfers \emph{representation dynamics} rather
than only representation states. We view normalized sentence embeddings as
points on a hypersphere and construct a teacher-conditioned semantic energy
that combines instance-level alignment with relational embedding geometry.
The negative Riemannian gradient of this energy, run under a finite-horizon
time reparameterization, defines a vector field that specifies a principled
direction of semantic evolution and reaches the teacher space exactly at the
final depth.
Instead of introducing a Neural ODE module or an external numerical solver,
we interpret the existing Transformer layers of the student as discrete
integration steps of this continuous flow and regularize consecutive layers
to follow the corresponding teacher-guided dynamics.

The resulting framework requires only cached final teacher embeddings,
does not access teacher hidden states during student training, and introduces
no additional module at inference time. We further establish that the
continuous dynamics preserve embedding norms and monotonically decrease the
teacher-conditioned distillation energy; in the instance-only case, the
geodesic distance to the teacher contracts in closed form along continuous
depth and vanishes exactly at the output layer.
We will evaluate whether this dynamics-based formulation improves
teacher-student transfer, representation geometry, and downstream
generalization over endpoint, multi-layer, relational, and flow-based
distillation baselines.
\end{abstract}


% ============================================================
% 1. INTRODUCTION
% ============================================================

\section{Introduction}\label{sec:introduction}

Dense text embeddings have become a fundamental interface between language
models and downstream applications, supporting semantic search,
retrieval-augmented generation, clustering, classification, and semantic
textual similarity. Recent embedding models obtain strong general-purpose
representations by increasing model size, training data, and representation
dimensionality. However, these improvements come with substantial inference
and deployment costs, motivating the compression of large embedding models
into smaller encoders through knowledge distillation
\citep{hinton2015distilling,reimers2020multilingual,xu2023distillcse,
zhang2024jasper,truong2025emo}.

A common principle underlying embedding distillation is
\emph{representation matching}. Given a teacher embedding
$\boldsymbol{\tau}$ and a student representation $\mathbf{z}$, the student
is optimized to minimize a discrepancy such as cosine distance or mean
squared error. More structured approaches enrich this supervision using
contrastive signals, intermediate hidden representations, attention
relations, optimal transport, or pairwise geometry
\citep{xu2023distillcse,truong2025emo}.
Recent work further explores depth-wise supervision. For example,
ESE encourages salient semantic information to emerge in shallower
representations to support depth-scalable sentence embeddings
\citep{li2025ese}, while TALAS anchors selected upper student layers to the
teacher output and propagates relational structure toward lower layers
through adjacent-layer self-distillation \citep{talas}.

Despite their differences, these approaches predominantly supervise a
collection of \emph{states}:
\begin{equation}
    \mathbf{z}^{(1)},\,
    \mathbf{z}^{(2)},\,
    \ldots,\,
    \mathbf{z}^{(L)}.
\end{equation}
The corresponding objectives describe what one or more intermediate
representations should resemble, but do not explicitly model the
\emph{transition} from one level of abstraction to the next.
We argue that this distinction matters for embedding distillation.
A shallow Transformer layer primarily captures local and lower-level
linguistic structure, whereas deeper layers gradually accumulate
task-relevant semantic information. Directly forcing many intermediate
layers toward a fixed high-level teacher representation therefore risks
destroying the student's natural hierarchical organization.
Indeed, recent layer-wise embedding distillation results indicate that
stronger supervision at an increasing number of student layers does not
necessarily produce better transfer, particularly under substantial
teacher-student capacity mismatch \citep{talas}.

This observation motivates a different question:

\begin{quote}
\emph{Instead of specifying where every student layer should be, can the
teacher specify how the student's representation should move through
depth?}
\end{quote}

The continuous-depth perspective of residual networks provides a natural
language for this question. Residual transformations can be interpreted as
discrete numerical approximations to an ordinary differential equation
(ODE), and this connection has motivated continuous-depth neural networks
and ODE-inspired Transformer architectures
\citep{chen2018neuralode,li2022odetransformer}.
Under this view, network depth plays the role of time and consecutive
hidden states form a discrete trajectory,
\begin{equation}
    \mathbf{z}^{(1)}
    \rightarrow
    \mathbf{z}^{(2)}
    \rightarrow \cdots \rightarrow
    \mathbf{z}^{(L)}.
\end{equation}
However, simply replacing Transformer layers with a generic Neural ODE does
not by itself solve the embedding distillation problem. Moreover,
continuous transformations have already been explored in knowledge
transfer through diffusion and flow-matching mechanisms
\citep{huang2023diffkd,shao2024fmkt}.
Our goal is therefore different: we do not learn an auxiliary continuous
network that transports a student output into a teacher output.
Instead, we use the teacher to define a \emph{semantic vector field}, and
train the student's existing discrete layers to realize that field.

We introduce \textbf{GeoODE-KD}, a teacher-guided geometric dynamics
framework for sentence embedding distillation. Our starting point is that
sentence embeddings are typically compared by cosine similarity and are
commonly $\ell_2$-normalized. We therefore model student and teacher
representations on the unit hypersphere
$\mathbb{S}^{d-1}$ rather than in unconstrained Euclidean space.
Given cached teacher embeddings, we construct a semantic potential energy
consisting of two complementary terms. The first attracts each student
sample toward its corresponding teacher representation, while the second
matches the pairwise geometry of the student batch to that of the teacher.
The negative Riemannian gradient of this energy defines a continuous vector
field,
\begin{equation}
    \frac{d\mathbf{Z}(t)}{dt}
    =
    -\frac{s(t)}{R(t)}
    \operatorname{grad}_{\mathbb{S}}
    \mathcal{E}\big(\mathbf{Z}(t),\mathbf{T}\big),
    \qquad
    R(t)=\int_t^1 s(u)\,du,
\end{equation}
which describes how sentence representations should evolve toward the
teacher's semantic space. The factor $1/R(t)$ is a time reparameterization
that turns the gradient flow, which would only reach its minimizer as
$t\to\infty$, into a finite-horizon flow that arrives at the teacher
exactly at $t=1$: in the instance-only case the trajectory is the geodesic
from the student to the teacher, traversed at the pace set by $s(t)$.

Crucially, GeoODE-KD does \emph{not} numerically integrate this equation
using an external ODE solver. We assign each Transformer layer a normalized
depth coordinate $t_\ell$ and treat consecutive student representations as
discrete observations of the continuous trajectory. At each depth, one
exact step of the teacher-conditioned flow, taken with the exponential map
of the hypersphere, provides the desired next representation. An ODE consistency objective then
aligns the actual next Transformer layer with this predicted state.
Consequently, the student itself becomes the numerical realization of the
teacher-guided semantic flow.

This construction has several desirable properties. First, the teacher is
used only to pre-compute final sentence embeddings; teacher inference and
teacher hidden states are unnecessary during student optimization. Second,
teacher supervision specifies a direction of semantic improvement rather
than forcing every intermediate layer to imitate the same terminal
representation. Third, because the dynamics live in the tangent space of
the hypersphere, they respect the geometry induced by normalized sentence
embeddings. Finally, all geometric and ODE-specific machinery is used only
as a training objective; deployment requires exactly the original student
encoder.

Our contributions are summarized as follows:
\begin{itemize}
    \item We formulate sentence embedding distillation as
    \emph{continuous representation dynamics}, shifting the target of
    knowledge transfer from isolated intermediate states to the evolution
    of representations through network depth.

    \item We introduce a teacher-conditioned Riemannian flow on the
    embedding hypersphere. Its energy jointly captures point-wise semantic
    alignment and teacher relational geometry, and a finite-horizon time
    reparameterization makes the flow reach the teacher exactly at unit
    depth, yielding an explicit and interpretable semantic vector field
    without a learnable Neural ODE module.

    \item We propose an ODE consistency objective that treats the
    Transformer's existing layers as discrete integration steps of the
    teacher-guided flow. The method requires only cached teacher output
    embeddings and incurs no additional inference-time architecture or ODE
    solver.

    \item We provide basic theoretical guarantees for the continuous
    dynamics, including hyperspherical norm preservation and monotonic
    distillation-energy descent. In the point-wise special case, we further
    derive the depth profile in closed form: the geodesic distance to the
    teacher contracts as $R(t)/R(0)$ and vanishes at the output layer, a
    quantitative prediction that the experiments can test layer by layer.
\end{itemize}


% ============================================================
% 2. RELATED WORK
% ============================================================

\section{Related Work}\label{sec:related}

\subsection{Knowledge Distillation for Sentence Embeddings}

Knowledge distillation transfers predictive or representational knowledge
from a high-capacity teacher to a compact student
\citep{hinton2015distilling}. In sentence representation learning,
\citet{reimers2020multilingual} demonstrated that sentence-level teacher
representations can directly supervise a student embedding space,
illustrating the effectiveness and simplicity of output-level embedding
distillation. Such methods are particularly attractive when the teacher is
large or accessible only through inference, since teacher embeddings can be
computed offline.

Subsequent approaches enrich the distillation signal beyond direct
point-wise imitation. DistillCSE incorporates knowledge distillation into
contrastive sentence representation learning and studies the instability
caused by high-variance teacher signals in contrastive settings
\citep{xu2023distillcse}. Jasper and Stella employ a multi-stage
distillation strategy for transferring strong embedding models into more
deployable representations \citep{zhang2024jasper}. More recently, EMO
addresses cross-tokenizer embedding distillation through intra-model
relational alignment and optimal transport over hidden representations
\citep{truong2025emo}.

Our work retains the efficiency advantage of sentence-level teacher
supervision but differs in what is distilled. Rather than using the final
teacher embedding solely as a static target, GeoODE-KD uses it to define a
vector field that supervises how student representations evolve through
depth.

\subsection{Layer-wise and Geometric Representation Transfer}

Intermediate supervision has long been used to expose students to richer
teacher information than final-output matching alone. In embedding models,
this idea is complicated by architectural and capacity differences:
representations at different depths need not encode semantic information at
comparable abstraction levels.

ESE addresses depth efficiency by explicitly encouraging salient
information to be represented at shallower layers, enabling sentence
embedding models to operate at different depths \citep{li2025ese}.
TALAS takes a complementary distillation perspective: selected upper
student layers are aligned to a cached final teacher embedding, while
lower-layer knowledge propagation is encouraged by matching pairwise
geometric relations between adjacent student layers \citep{talas}.
Its formulation highlights an important tension: direct teacher alignment
can provide high-level semantics, but excessive anchoring may over-constrain
lower or intermediate student representations.

GeoODE-KD builds on the broader insight that embedding geometry contains
knowledge not captured by point-wise targets. However, rather than imposing
a sequence of static layer constraints or requiring each adjacent layer to
preserve the same relation matrix, we define a \emph{global potential}
conditioned on the teacher geometry. Its gradient determines the desired
local transition at every student depth. Thus, relational knowledge affects
not only the location of embeddings but also their direction of evolution.

\subsection{Continuous-Depth Networks and ODE-Based Knowledge Transfer}

Neural ODEs interpret hidden representations as states of a continuous
dynamical system,
\begin{equation}
    \frac{d\mathbf{h}(t)}{dt}
    =
    f_\theta(\mathbf{h}(t),t),
\end{equation}
whose solution replaces an explicitly specified sequence of discrete
layers \citep{chen2018neuralode}. The connection between residual networks
and numerical ODE integration has also been explored for Transformers.
ODE Transformer, for example, relates Transformer residual computation to
numerical integration and derives architectures inspired by higher-order
ODE solvers \citep{li2022odetransformer}.

Continuous transformations have separately appeared in knowledge
distillation. DiffKD models the student feature as a noisy version of the
teacher feature and uses diffusion-based denoising before feature matching
\citep{huang2023diffkd}. Flow Matching provides a simulation-free framework
for learning continuous normalizing flows \citep{lipman2023flow}, extended
to Riemannian manifolds by Riemannian Flow Matching, whose reference
trajectories are geodesics traversed at a prescribed schedule
\citep{chen2024riemannian}, and FM-KT applies continuous flow-based
transport to progressively bridge teacher-student feature or logit
distributions \citep{shao2024fmkt}. Our finite-horizon field shares the
geodesic reference of Riemannian Flow Matching, but no velocity network is
learned: the student's own layers are trained to realize the analytic flow.

GeoODE-KD is conceptually distinct from both continuous-depth architecture
design and learned flow-based feature transport. We neither replace the
student architecture with a Neural ODE block nor train an auxiliary
continuous normalizing flow that maps a student output to a teacher output.
The student Transformer remains unchanged. Continuous dynamics instead
serve as a \emph{training principle}: the teacher induces an analytic
geometric vector field, and ordinary student layers are optimized to
approximate its discretized trajectory. As a result, no ODE solver or
flow module is required at inference time.


% ============================================================
% 3. METHOD
% ============================================================

\section{GeoODE-KD}\label{sec:method}

\subsection{Problem Formulation}\label{sec:problem}

Let
\begin{equation}
    f_T : \mathcal{X}\rightarrow\mathbb{R}^{d_T},
    \qquad
    f_S : \mathcal{X}\rightarrow\mathbb{R}^{d_S}
\end{equation}
denote a teacher embedding model and a compact student encoder,
respectively. Given an unlabeled distillation corpus
\begin{equation}
    \mathcal{D}={\bigl\{x_i\bigr\}}_{i=1}^{N},
\end{equation}
the objective is to transfer the semantic structure of $f_T$ into $f_S$
without requiring task labels.

The student contains $L$ Transformer layers. For sentence $x_i$, let
$\mathbf{h}^{(\ell)}_i$ denote its hidden representation at layer $\ell$.
A pooling operator $\operatorname{Pool}(\cdot)$ produces an intermediate
sentence representation,
\begin{equation}
    \mathbf{q}^{(\ell)}_i
    =
    \operatorname{Pool}
    \left(
        \mathbf{h}^{(\ell)}_i
    \right).
\end{equation}

We focus on the common large-teacher-small-student regime where
$d_T\ge d_S$. Teacher outputs are transformed once into the student
dimension using a fixed linear map
$P_T=P_{\mathrm{PCA}}R\in\mathbb{R}^{d_T\times d_S}$, where
$P_{\mathrm{PCA}}$ is PCA fitted only on the cached training embeddings and
$R\in O(d_S)$ is a fixed rotation inside the PCA subspace whose role is
explained below:
\begin{equation}
    \boldsymbol{\tau}_i
    =
    \operatorname{norm}
    \left(
        f_T(x_i)P_{\mathrm{PCA}}R
    \right),
    \label{eq:teacher_embedding}
\end{equation}
where
\begin{equation}
    \operatorname{norm}(\mathbf{x})
    =
    \frac{\mathbf{x}}
         {\|\mathbf{x}\|_2+\epsilon}.
\end{equation}
When teacher and student dimensions are equal, $P_T$ is the identity.
PCA is not merely a convenient choice here: by the Eckart--Young theorem it
is, among all rank-$d_S$ linear maps, the one that retains the largest share
of the embedding energy and therefore best preserves the Gram matrix
$\mathbf{T}\mathbf{T}^{\top}$ on which the relational energy of
Section~\ref{sec:energy} is defined. In our main setting
($d_T=2560\rightarrow d_S=768$) the fitted map retains $93.0\%$ of the
cached embedding energy.

\paragraph{Gauge freedom of the endpoint.}
Every quantity by which a sentence encoder is evaluated, cosine
similarity, Spearman correlation, linear probes, is invariant under an
orthogonal transformation $\mathbf{z}\mapsto\mathbf{z}R$ of the student
space, and so is the relational energy of Section~\ref{sec:energy}. The
coordinates that PCA returns are therefore an arbitrary \emph{gauge}: any
rotation of them inside the retained subspace is an equally valid target,
but a point-wise objective such as Eq.~\eqref{eq:semantic_energy} pins one
of them. Left at the PCA gauge, the endpoint condition asks the student to
rotate its entire pretrained representation space onto the teacher's
principal axes, a demand unrelated to semantics: in our setting the
untrained student is nearly orthogonal to such targets (mean cosine
$\approx0$). We remove the gauge by fixing $R$ once, before training, as
the orthogonal Procrustes alignment \citep{schonemann1966procrustes} of the
PCA coordinates to the untrained student,
\begin{equation}
    R
    =
    \underset{R\in O(d_S)}{\arg\min}
    \left\|
        \mathbf{T}_{\mathrm{PCA}}R-\mathbf{Z}^{(L)}_{0}
    \right\|_F
    =
    UV^{\top},
    \qquad
    U\Sigma V^{\top}
    =
    \operatorname{svd}\!\left(
        \mathbf{T}_{\mathrm{PCA}}^{\top}\mathbf{Z}^{(L)}_{0}
    \right),
    \label{eq:gauge}
\end{equation}
where $\mathbf{T}_{\mathrm{PCA}}$ stacks the PCA coordinates of a subset of
the corpus and $\mathbf{Z}^{(L)}_{0}$ the initial student's output
embeddings of the same sentences. Since $R$ is orthogonal, the Gram matrix
$\mathbf{T}\mathbf{T}^{\top}$, the retained variance and the
Eckart--Young optimality of $P_{\mathrm{PCA}}$ are untouched; only the
arbitrary part of the endpoint condition is removed. Like
$P_{\mathrm{PCA}}$, $R$ is a closed-form statistic of the cached targets
and the initial student, computed once (one SVD of a
$d_S\times d_S$ matrix), frozen, and discarded at inference. $P_T$ never
interacts with student optimization.

Student representations are similarly normalized,
\begin{equation}
    \mathbf{z}^{(\ell)}_i
    =
    \operatorname{norm}
    \left(
        \mathbf{q}^{(\ell)}_i
    \right).
    \label{eq:student_norm}
\end{equation}

Therefore,
\begin{equation}
    \mathbf{z}^{(\ell)}_i,\,
    \boldsymbol{\tau}_i
    \in
    \mathbb{S}^{d_S-1},
\end{equation}
where $\mathbb{S}^{d_S-1}$ denotes the unit hypersphere.

For a mini-batch of size $B$, we stack row-wise representations as
\begin{equation}
    \mathbf{Z}^{(\ell)}
    =
    \begin{bmatrix}
        {\bigl(\mathbf{z}^{(\ell)}_1\bigr)}^\top\\
        \vdots\\
        {\bigl(\mathbf{z}^{(\ell)}_B\bigr)}^\top
    \end{bmatrix}
    \in\mathbb{R}^{B\times d_S},
    \qquad
    \mathbf{T}
    =
    \begin{bmatrix}
        \boldsymbol{\tau}_1^\top\\
        \vdots\\
        \boldsymbol{\tau}_B^\top
    \end{bmatrix}.
\end{equation}


% ------------------------------------------------------------
\subsection{From Discrete Layer Matching to Continuous Semantic Dynamics}\label{sec:continuous}

Conventional multi-layer distillation can be written schematically as
\begin{equation}
    \mathcal{L}_{\mathrm{layer}}
    =
    \sum_{\ell\in\mathcal{A}}
    D\left(
        \mathbf{Z}^{(\ell)},
        \mathbf{T}
    \right),
    \label{eq:static_layer}
\end{equation}
where $\mathcal{A}$ denotes a selected set of student layers.
Equation~\eqref{eq:static_layer} treats each supervised layer as an
independent state-matching problem.

We instead introduce a normalized depth coordinate
\begin{equation}
    t_\ell=\frac{\ell}{L},
    \qquad
    \ell\in\{1,\ldots,L\},
\end{equation}
and interpret
\begin{equation}
    \mathbf{Z}^{(\ell)}
    \approx
    \mathbf{Z}(t_\ell)
\end{equation}
as samples from an underlying continuous semantic trajectory.

Our objective is to construct a teacher-conditioned differential equation
\begin{equation}
    \frac{d\mathbf{Z}(t)}{dt}
    =
    F\left(
        \mathbf{Z}(t),\mathbf{T},t
    \right)
    \label{eq:general_ode}
\end{equation}
and train the discrete Transformer layers to approximate this flow.


% ------------------------------------------------------------
\subsection{Teacher-Conditioned Semantic Energy}\label{sec:energy}

The vector field in Eq.~\eqref{eq:general_ode} should satisfy two
properties. First, each sentence representation should progressively move
toward the semantic target supplied by its teacher embedding. Second, the
student should recover the relational geometry of the teacher embedding
space rather than merely matching samples independently.

We encode these requirements through a teacher-conditioned energy
function.

\paragraph{Instance-level semantic energy.}

For two points on the unit hypersphere, the intrinsic notion of discrepancy
is the geodesic (great-circle) distance
$d_g(\mathbf{z},\boldsymbol{\tau})=\arccos(\mathbf{z}^{\top}\boldsymbol{\tau})$.
We define
\begin{equation}
    \mathcal{E}_{\mathrm{sem}}(\mathbf{Z},\mathbf{T})
    =
    \frac{1}{2B}
    \sum_{i=1}^{B}
    d_g^2\left(\mathbf{z}_i,\boldsymbol{\tau}_i\right).
    \label{eq:semantic_energy}
\end{equation}

The squared geodesic distance is the natural potential on a Riemannian
manifold: its negative Riemannian gradient is the logarithmic map
$\operatorname{Log}_{\mathbf{z}}(\boldsymbol{\tau})$, the tangent vector
that points along the geodesic to $\boldsymbol{\tau}$ and whose length is
$d_g$. Near the teacher it agrees with the cosine distance,
$\tfrac12 d_g^2 = 1-\mathbf{z}^{\top}\boldsymbol{\tau}+O(d_g^4)$, but its
gradient does not vanish for far-away states, which is what will allow the
flow below to reach the teacher in finite depth.

\paragraph{Relational geometric energy.}

Sentence embeddings are primarily useful through relative similarities.
We therefore construct batch-wise Gram matrices
\begin{equation}
    \mathbf{G}_{S}
    =
    \mathbf{Z}\mathbf{Z}^{\top},
    \qquad
    \mathbf{G}_{T}
    =
    \mathbf{T}\mathbf{T}^{\top}.
\end{equation}
Their $(i,j)$-th entries correspond to pairwise cosine similarities.

The geometric energy is
\begin{equation}
    \mathcal{E}_{\mathrm{geo}}(\mathbf{Z},\mathbf{T})
    =
    \frac{1}{B^2}
    \left\|
        \mathbf{Z}\mathbf{Z}^{\top}
        -
        \mathbf{T}\mathbf{T}^{\top}
    \right\|_F^2.
    \label{eq:geometry_energy}
\end{equation}

The total teacher-conditioned potential is
\begin{equation}
    \boxed{
    \mathcal{E}(\mathbf{Z},\mathbf{T})
    =
    \alpha
    \mathcal{E}_{\mathrm{sem}}(\mathbf{Z},\mathbf{T})
    +
    \beta
    \mathcal{E}_{\mathrm{geo}}(\mathbf{Z},\mathbf{T})
    }
    \label{eq:total_energy}
\end{equation}
with $\alpha,\beta\ge0$.

Importantly, Eq.~\eqref{eq:total_energy} is not directly minimized at
every student layer. Instead, it defines the continuous dynamics that
student layers are trained to realize.


% ------------------------------------------------------------
\subsection{Teacher-Guided Finite-Horizon Riemannian Flow}\label{sec:riemannian_flow}

Because the embeddings lie on
$\mathbb{S}^{d_S-1}$, unconstrained Euclidean motion is undesirable:
an arbitrary update can alter vector norms and introduce radial components
that are irrelevant to cosine geometry. We therefore construct a
Riemannian flow on the product of hyperspheres, driven by the negative
gradient of $\mathcal{E}$ and reparameterized in time so that it reaches
the teacher within the unit depth interval.

For an embedding $\mathbf{z}$, the tangent space is
\begin{equation}
    T_{\mathbf{z}}\mathbb{S}^{d_S-1}
    =
    \{
        \mathbf{v}\in\mathbb{R}^{d_S}
        :
        \mathbf{z}^{\top}\mathbf{v}=0
    \}.
\end{equation}

The orthogonal projection of a Euclidean vector $\mathbf{u}$ onto this
tangent space is
\begin{equation}
    \Pi_{\mathbf{z}}(\mathbf{u})
    =
    \mathbf{u}
    -
    (\mathbf{z}^{\top}\mathbf{u})
    \mathbf{z}.
    \label{eq:tangent_projection}
\end{equation}

For the batch representation, $\Pi_{\mathbf{Z}}$ applies
Eq.~\eqref{eq:tangent_projection} row-wise.

The negative Riemannian gradient of the semantic energy is the row-wise
logarithmic map of the hypersphere,
\begin{equation}
    -\operatorname{grad}_{\mathbb{S}}
    \mathcal{E}_{\mathrm{sem}}
    =
    \frac{1}{B}\operatorname{Log}_{\mathbf{Z}}(\mathbf{T}),
    \qquad
    \operatorname{Log}_{\mathbf{z}}(\boldsymbol{\tau})
    =
    \frac{d_g(\mathbf{z},\boldsymbol{\tau})}
         {\sin d_g(\mathbf{z},\boldsymbol{\tau})}
    \Pi_{\mathbf{z}}(\boldsymbol{\tau}),
    \label{eq:log_map}
\end{equation}
which is already tangent to the sphere, whereas the Euclidean gradient of
the geometric term satisfies
\begin{equation}
    \nabla_{\mathbf{Z}}
    \mathcal{E}_{\mathrm{geo}}
    =
    \frac{4}{B^2}
    \left(
        \mathbf{Z}\mathbf{Z}^{\top}
        -
        \mathbf{T}\mathbf{T}^{\top}
    \right)\mathbf{Z}.
    \label{eq:geo_gradient}
\end{equation}

We therefore define the tangent descent direction
\begin{align}
    \mathbf{V}(\mathbf{Z},\mathbf{T})
    &=
    \alpha\operatorname{Log}_{\mathbf{Z}}(\mathbf{T})
    -
    \frac{4\beta}{B}
    \Pi_{\mathbf{Z}}
    \left[
    \left(
        \mathbf{Z}\mathbf{Z}^{\top}
        -
        \mathbf{T}\mathbf{T}^{\top}
    \right)\mathbf{Z}
    \right]
    =
    -B\operatorname{grad}_{\mathbb{S}}\mathcal{E}(\mathbf{Z},\mathbf{T}),
    \label{eq:euclidean_field}\\
    F(\mathbf{Z},\mathbf{T},t)
    &=
    \frac{s(t)}{R(t)}
    \mathbf{V}(\mathbf{Z},\mathbf{T}),
    \qquad
    R(t)=\int_t^1 s(u)\,du .
    \label{eq:vector_field}
\end{align}
The factor $B$ in Eq.~\eqref{eq:euclidean_field} removes the $1/B$ that
differentiating a batch mean introduces, so that the speed of the flow does
not depend on the batch size; it does not change the direction.

The continuous teacher-guided dynamics are then
\begin{equation}
    \boxed{
    \frac{d\mathbf{Z}(t)}{dt}
    =
    F(\mathbf{Z}(t),\mathbf{T},t)
    =
    -
    \frac{s(t)}{R(t)}
    \,B\operatorname{grad}_{\mathbb{S}}
    \mathcal{E}
    \left(
        \mathbf{Z}(t),\mathbf{T}
    \right)
    }
    \label{eq:geoode}
\end{equation}

where $s(t)\ge0$ controls how teacher guidance is distributed over depth
and $R(t)$ is the guidance that remains to be spent after depth $t$.

\paragraph{Why a finite-horizon reparameterization.}
A plain gradient flow $\dot{\mathbf{Z}}=-s(t)\operatorname{grad}\mathcal{E}$
converges to the teacher only as $t\to\infty$, because the gradient
vanishes at the minimizer. Over the unit depth interval it can cover at most
$\int_0^1 s(t)\,\|\mathbf{V}\|\,dt$ of geodesic length, whereas a student
whose first layer is nearly orthogonal to the teacher has to travel about
$\pi/2$. With $s(t)=t$ this is at most $0.5$ against $1.6$ radians, so such
a flow cannot reach the teacher, and a consistency loss built on it can only
be satisfied by moving slowly at every depth and leaving the output layer to
make the jump alone. The reparameterization by $1/R(t)$ is the
substitution $d\tau=s(t)\,dt/R(t)$, under which the gradient flow is run
for infinite $\tau$-time while $t$ only spans $[0,1]$. Corollary~1 below
shows that the resulting instance-only trajectory is exactly the geodesic
to the teacher, traversed so that the remaining distance is proportional to
the remaining guidance $R(t)$.

\paragraph{Depth-dependent semantic guidance.}

We use
\begin{equation}
    s(t)=t
    \label{eq:schedule}
\end{equation}
as the default schedule.
This choice introduces weak teacher guidance at shallow depths and
progressively strengthens it toward the output layers. The design reflects
the intuition that shallow representations should retain flexibility to
extract lower-level linguistic information, while deeper layers should
become increasingly organized by the teacher's semantic geometry.

With this choice $R(t)=(1-t^2)/2$, so the fraction of the remaining
geodesic covered between consecutive depths grows with depth (Eq.~\eqref{eq:step_fraction}).
More general schedules,
e.g.,
$s(t)=t^p$ with $R(t)=(1-t^{p+1})/(p+1)$, or a constant $s(t)=1$ with
$R(t)=1-t$, can be considered as ablations; none introduces learned
parameters, and all of them reach the teacher at $t=1$.


% ------------------------------------------------------------
\subsection{Student Layers as Discrete ODE Steps}\label{sec:discrete}

A central design choice of GeoODE-KD is that
Eq.~\eqref{eq:geoode} is \emph{not} solved using a black-box ODE solver.
Instead, the student Transformer itself provides the discrete trajectory.

For consecutive student depths,
\begin{equation}
    \Delta t
    =
    t_{\ell+1}-t_\ell
    =
    \frac{1}{L}.
\end{equation}

The time warp in Eq.~\eqref{eq:vector_field} is integrated exactly over
the layer interval rather than evaluated at its left endpoint. Since
$\dot R=-s$, the warped time elapsed between two depths is
$\int_{t_\ell}^{t_{\ell+1}} s(u)/R(u)\,du=\log\bigl(R(t_\ell)/R(t_{\ell+1})\bigr)$,
and the corresponding step fraction is
\begin{equation}
    \rho_\ell
    =
    1-\exp\left(-\int_{t_\ell}^{t_{\ell+1}}\frac{s(u)}{R(u)}\,du\right)
    =
    1-\frac{R(t_{\ell+1})}{R(t_\ell)}
    \in(0,1],
    \label{eq:step_fraction}
\end{equation}
which is the fraction of the remaining geodesic to the target that the flow
covers between depths $t_\ell$ and $t_{\ell+1}$. For $s(t)=t$,
$\rho_\ell=(2\ell+1)/(L^2-\ell^2)$, and $\rho_{L-1}=1$ for every
schedule.

Given $\mathbf{Z}^{(\ell)}$, one step of the teacher-guided dynamics is
\begin{equation}
    \widetilde{\mathbf{Z}}^{(\ell+1)}
    =
    \operatorname{Exp}_{\mathbf{Z}^{(\ell)}}
    \left(
        \rho_\ell\,
        \mathbf{V}(
            \mathbf{Z}^{(\ell)},
            \mathbf{T}
        )
    \right),
    \label{eq:riemannian_euler}
\end{equation}
where $\operatorname{Exp}$ is the exponential map of the hypersphere,
applied row-wise and available in closed form:
\begin{equation}
    \operatorname{Exp}_{\mathbf{z}}(\mathbf{v})
    =
    \cos\left(\|\mathbf{v}\|\right)\mathbf{z}
    +
    \sin\left(\|\mathbf{v}\|\right)
    \frac{\mathbf{v}}{\|\mathbf{v}\|}.
    \label{eq:retraction}
\end{equation}
Unlike a first-order retraction such as row-wise normalization of
$\mathbf{Z}+\mathbf{v}$, which rotates by $\arctan\|\mathbf{v}\|$
instead of $\|\mathbf{v}\|$ and would fall short on the large final step,
the exponential map is exact for every step size. In the instance-only case
($\beta=0$), Eq.~\eqref{eq:riemannian_euler} is precisely spherical linear
interpolation, $\widetilde{\mathbf{z}}^{(\ell+1)}=\operatorname{slerp}(\mathbf{z}^{(\ell)},\boldsymbol{\tau};\rho_\ell)$,
so the predicted state for the output layer is the teacher itself.

The actual Transformer produces $\mathbf{Z}^{(\ell+1)}$.
We regularize it toward the ODE-predicted next state:
\begin{equation}
    \mathcal{L}_{\mathrm{dyn}}
    =
    \frac{1}{L-1}
    \sum_{\ell=1}^{L-1}
    D_{\cos}
    \left(
        \mathbf{Z}^{(\ell+1)},
        \operatorname{sg}
        \left[
            \widetilde{\mathbf{Z}}^{(\ell+1)}
        \right]
    \right),
    \label{eq:dynamics_loss}
\end{equation}
where $\operatorname{sg}[\cdot]$ denotes stop-gradient and
\begin{equation}
    D_{\cos}(\mathbf{A},\mathbf{B})
    =
    \frac{1}{B}
    \sum_{i=1}^{B}
    \left(
        1-
        \mathbf{a}_i^\top\mathbf{b}_i
    \right).
\end{equation}

The stop-gradient operation makes
$\widetilde{\mathbf{Z}}^{(\ell+1)}$ a local dynamics target computed from
the current trajectory, preventing the model from trivially modifying both
sides of the consistency objective simultaneously. We will separately
ablate full-gradient dynamics.

Equation~\eqref{eq:dynamics_loss} creates the main distinction between
GeoODE-KD and static multi-layer distillation:
\begin{align}
    \text{State matching:}
    &&
    \mathbf{Z}^{(\ell)}
    &\rightarrow
    \mathbf{T},
    \\
    \text{GeoODE-KD:}
    &&
    \mathbf{Z}^{(\ell)}
    &\rightarrow
    \mathbf{Z}^{(\ell+1)}
    \quad
    \text{according to }
    F(
        \mathbf{Z}^{(\ell)},\mathbf{T},t_\ell
    ).
\end{align}

The teacher thus supervises the student's \emph{velocity through depth}
instead of demanding that every intermediate representation imitate the
same terminal target: each layer is asked to cover a prescribed fraction
$\rho_\ell$ of what remains, and the fractions compose to the whole
geodesic.


% ------------------------------------------------------------
\subsection{Endpoint Semantic Distillation}\label{sec:endpoint}

While $\mathcal{L}_{\mathrm{dyn}}$ determines the local trajectory, the
student output should ultimately remain close to the teacher endpoint.
We therefore apply direct distillation only at the final student layer:
\begin{equation}
    \boxed{
    \mathcal{L}_{\mathrm{end}}
    =
    \frac{1}{B}
    \sum_{i=1}^{B}
    \left(
        1-
        {\bigl(\mathbf{z}^{(L)}_i\bigr)}^\top
        \boldsymbol{\tau}_i
    \right)
    }.
    \label{eq:endpoint_loss}
\end{equation}

Unlike conventional multi-layer teacher anchoring, no intermediate student
layer is directly forced to equal the final teacher embedding. Intermediate
supervision occurs only through the vector field in
Eq.~\eqref{eq:vector_field}. Because $\rho_{L-1}=1$, the dynamics target
for the output layer already coincides with $\boldsymbol{\tau}_i$ when
$\beta=0$; $\mathcal{L}_{\mathrm{end}}$ is therefore the boundary
condition of the same flow rather than a competing objective, and it
remains the only term that anchors the output when the relational term
bends the final step.


% ------------------------------------------------------------
\subsection{Contrastive Representation Regularization}\label{sec:contrastive}

Pure teacher imitation may transfer undesirable concentration or
anisotropy from the teacher embedding distribution. We therefore retain a
standard unsupervised contrastive objective at the final layer.

For each input $x_i$, two independent dropout masks produce normalized
representations
$\mathbf{z}^{(L,a)}_i$ and $\mathbf{z}^{(L,b)}_i$.
We use
\begin{equation}
    \mathcal{L}_{\mathrm{ctr}}
    =
    -
    \frac{1}{B}
    \sum_{i=1}^{B}
    \log
    \frac{
        \exp
        \left(
            \operatorname{sim}
            (
                \mathbf{z}^{(L,a)}_i,
                \mathbf{z}^{(L,b)}_i
            )/\tau_c
        \right)
    }{
        \sum_{j=1}^{B}
        \exp
        \left(
            \operatorname{sim}
            (
                \mathbf{z}^{(L,a)}_i,
                \mathbf{z}^{(L,b)}_j
            )/\tau_c
        \right)
    },
    \label{eq:contrastive}
\end{equation}
where $\tau_c$ is the contrastive temperature.

Importantly, this objective is applied only to the final representation.
We avoid imposing contrastive targets independently at every layer, which
would confound the role of continuous dynamics.


% ------------------------------------------------------------
\subsection{Overall Training Objective}\label{sec:objective}

The final GeoODE-KD objective is
\begin{equation}
    \boxed{
    \mathcal{L}_{\mathrm{total}}
    =
    \lambda_{\mathrm{end}}
    \mathcal{L}_{\mathrm{end}}
    +
    \lambda_{\mathrm{dyn}}
    \mathcal{L}_{\mathrm{dyn}}
    +
    \lambda_{\mathrm{ctr}}
    \mathcal{L}_{\mathrm{ctr}}
    }.
    \label{eq:total_loss}
\end{equation}

The method introduces two conceptually distinct hyperparameters inside the
semantic potential:
$\alpha$ controls point-wise teacher attraction and $\beta$ controls
relational geometric preservation.
For the initial implementation, we set $\alpha=1$ and tune only $\beta$,
while fixing
$\lambda_{\mathrm{end}}=\lambda_{\mathrm{dyn}}=1$ and tuning
$\lambda_{\mathrm{ctr}}$ on a validation split.

We emphasize the separation between
Eq.~\eqref{eq:total_energy} and Eq.~\eqref{eq:total_loss}.
The energy
$\mathcal{E}$ is not simply another layer-wise training loss.
Instead, it defines the vector field that specifies the desired
\emph{direction of change}, whereas
$\mathcal{L}_{\mathrm{dyn}}$ measures whether the Transformer's actual
layer transition follows that direction.


% ------------------------------------------------------------
\subsection{Theoretical Properties}\label{sec:theory}

We next state basic properties that motivate the proposed dynamics.

\paragraph{Proposition 1: Hyperspherical norm preservation.}

Consider a single trajectory
$\mathbf{z}(t)\in\mathbb{S}^{d-1}$ governed by
\begin{equation}
    \dot{\mathbf{z}}
    =
    \Pi_{\mathbf{z}}(\mathbf{u}).
\end{equation}
Then
\begin{equation}
    \frac{d}{dt}
    \|\mathbf{z}(t)\|_2^2
    =
    0.
\end{equation}

\paragraph{Proof.}
By construction,
$\Pi_{\mathbf{z}}(\mathbf{u})$ lies in the tangent space of the
hypersphere and therefore
\begin{equation}
    \mathbf{z}^{\top}
    \Pi_{\mathbf{z}}(\mathbf{u})
    =
    0.
\end{equation}
Consequently,
\begin{equation}
    \frac{d}{dt}
    \|\mathbf{z}\|_2^2
    =
    2\mathbf{z}^{\top}\dot{\mathbf{z}}
    =
    0.
    \qquad \square
\end{equation}

Thus, the continuous teacher-guided dynamics alter semantic direction
without introducing an irrelevant radial component.


\paragraph{Proposition 2: Monotonic teacher-conditioned energy descent.}

Let the continuous trajectory follow Eq.~\eqref{eq:geoode},
\begin{equation}
    \dot{\mathbf{Z}}
    =
    -
    \frac{s(t)}{R(t)}
    \,B\operatorname{grad}_{\mathbb{S}}
    \mathcal{E}(\mathbf{Z},\mathbf{T}),
\end{equation}
where $s(t)\ge0$ and $t<1$. Then
\begin{equation}
    \frac{d}{dt}
    \mathcal{E}(\mathbf{Z}(t),\mathbf{T})
    \le0.
\end{equation}

\paragraph{Proof.}
Using the Riemannian chain rule,
\begin{align}
    \frac{d\mathcal{E}}{dt}
    &=
    \left\langle
        \operatorname{grad}_{\mathbb{S}}\mathcal{E},
        \dot{\mathbf{Z}}
    \right\rangle
    \\
    &=
    -
    \frac{B\,s(t)}{R(t)}
    \left\|
        \operatorname{grad}_{\mathbb{S}}\mathcal{E}
    \right\|_F^2
    \le0.
    \qquad \square
\end{align}

Hence, the ideal continuous trajectory monotonically reduces the combined
semantic and relational discrepancy from the teacher.


\paragraph{Corollary 1: Geodesic contraction and exact arrival under the point-wise flow.}

Consider $\beta=0$, $\alpha=1$, and a single normalized teacher-student pair
$(\boldsymbol{\tau},\mathbf{z})$. The dynamics reduce to
\begin{equation}
    \dot{\mathbf{z}}
    =
    \frac{s(t)}{R(t)}
    \operatorname{Log}_{\mathbf{z}}(\boldsymbol{\tau}).
    \label{eq:anchor_flow}
\end{equation}
Let $\theta(t)=d_g(\mathbf{z}(t),\boldsymbol{\tau})$. Since
$\operatorname{Log}_{\mathbf{z}}(\boldsymbol{\tau})$ is the unit-speed
geodesic direction scaled by $\theta$, the trajectory stays on the geodesic
from $\mathbf{z}(0)$ to $\boldsymbol{\tau}$ and
\begin{equation}
    \dot\theta(t)
    =
    -\frac{s(t)}{R(t)}\,\theta(t),
    \qquad
    \dot R(t)=-s(t)
    \quad\Longrightarrow\quad
    \frac{d}{dt}\log\theta
    =
    \frac{d}{dt}\log R .
\end{equation}
Hence
\begin{equation}
    \boxed{
    \theta(t)
    =
    \theta(0)\,\frac{R(t)}{R(0)}
    }
    \qquad\text{and}\qquad
    \theta(1)=0 .
    \label{eq:monotonic_cosine}
\end{equation}

The geodesic distance to the teacher therefore decreases monotonically,
in proportion to the guidance that remains, and vanishes exactly at unit
depth: with the default $s(t)=t$, $\theta(t)=\theta(0)(1-t^2)$. Cosine
similarity $\cos\theta(t)$ increases monotonically as a consequence.
Equation~\eqref{eq:monotonic_cosine} is a quantitative prediction for the
depth profile of a student that realizes the flow, and the discrete steps
of Eq.~\eqref{eq:riemannian_euler} reproduce it exactly at the sampled
depths: $\theta(t_{\ell+1})=(1-\rho_\ell)\theta(t_\ell)$.
When relational geometry is included, individual distances need not
follow Eq.~\eqref{eq:monotonic_cosine}, but Proposition~2 guarantees
monotonic descent of the joint semantic-geometric energy.


% ------------------------------------------------------------
\subsection{Training Algorithm}\label{sec:algorithm}

\begin{algorithm}[t]
\caption{GeoODE-KD Training}\label{alg:geood}
\begin{algorithmic}[1]
\Require{}
Unlabeled corpus $\mathcal{D}$,
teacher $f_T$,
student $f_S$ with $L$ layers,
$\alpha,\beta$,
$\lambda_{\mathrm{end}},
\lambda_{\mathrm{dyn}},
\lambda_{\mathrm{ctr}}$

\State{} Run $f_T$ once over $\mathcal{D}$ and cache teacher embeddings
\State{} Fit fixed PCA map $P_{\mathrm{PCA}}$ on cached teacher embeddings
\State{} Fix the gauge $R$ by Procrustes against the initial student, Eq.~\eqref{eq:gauge}
\State{} Transform and normalize all cached targets using
Eq.~\eqref{eq:teacher_embedding}

\For{each training mini-batch $\mathcal{B}$}
    \State{} Retrieve cached teacher matrix $\mathbf{T}$

    \State{} Run the student and obtain
    ${\bigl\{\mathbf{Z}^{(\ell)}\bigr\}}_{\ell=1}^{L}$

    \State{} $\mathcal{L}_{\mathrm{dyn}}\gets0$

    \For{$\ell=1,\ldots,L-1$}
        \State{} $t_\ell\gets \ell/L$

        \State{} Compute
        $\mathbf{G}_S
        \gets
        \mathbf{Z}^{(\ell)}
        {\bigl(\mathbf{Z}^{(\ell)}\bigr)}^\top$

        \State{} Compute
        $\mathbf{G}_T
        \gets
        \mathbf{T}\mathbf{T}^\top$

        \State{} Compute the tangent direction
        $\mathbf{V}(
            \mathbf{Z}^{(\ell)},
            \mathbf{T}
        )$
        using Eq.~\eqref{eq:euclidean_field}

        \State{} $\rho_\ell\gets 1-R(t_{\ell+1})/R(t_\ell)$
        using Eq.~\eqref{eq:step_fraction}

        \State{} Predict
        $\widetilde{\mathbf{Z}}^{(\ell+1)}
        \gets
        \operatorname{Exp}_{\mathbf{Z}^{(\ell)}}(\rho_\ell\mathbf{V})$
        using Eq.~\eqref{eq:riemannian_euler}

        \State{} Accumulate ODE consistency using
        Eq.~\eqref{eq:dynamics_loss}
    \EndFor{}

    \State{} Compute endpoint loss
    $\mathcal{L}_{\mathrm{end}}$
    using Eq.~\eqref{eq:endpoint_loss}

    \State{} Obtain a second dropout view and compute
    $\mathcal{L}_{\mathrm{ctr}}$

    \State{} Compute
    $\mathcal{L}_{\mathrm{total}}$
    using Eq.~\eqref{eq:total_loss}

    \State{} Update student parameters with AdamW
\EndFor{}
\end{algorithmic}
\end{algorithm}


% ------------------------------------------------------------
\subsection{Training and Inference Complexity}\label{sec:complexity}

Teacher inference is entirely offline. For a corpus of $N$ examples, only
one final teacher embedding per training example is stored. During student
training, GeoODE-KD does not require teacher hidden states, attention maps,
or parallel teacher forward passes.

For a batch of $B$ samples and embedding dimension $d_S$, computing the
relational component at one layer requires
$\mathcal{O}(B^2d_S)$ operations for the Gram matrix and
$\mathcal{O}(B^2)$ temporary relation memory. Across $L$ student layers,
the dynamics objective therefore adds approximately
\begin{equation}
    \mathcal{O}(LB^2d_S)
\end{equation}
training computation. This cost is independent of the size of the teacher.

At inference time, all teacher-conditioned components are removed:
\begin{equation}
    x
    \xrightarrow[]{f_S}
    \mathbf{z}^{(L)}.
\end{equation}
No PCA target projection, gauge rotation, Gram matrix, semantic energy, vector field,
ODE solver, or auxiliary network is evaluated. GeoODE-KD therefore has
the same inference architecture and asymptotic inference cost as its
student backbone.


% ------------------------------------------------------------
\subsection{Relationship to Alternative Distillation Objectives}\label{sec:comparison}

GeoODE-KD can be understood by contrasting four forms of supervision.

\paragraph{Endpoint KD.}
\begin{equation}
    \mathbf{Z}^{(L)}
    \rightarrow
    \mathbf{T}.
\end{equation}
Only the final state is constrained.

\paragraph{Static multi-layer KD.}
\begin{equation}
    \mathbf{Z}^{(\ell)}
    \rightarrow
    \mathbf{T}
    \qquad
    \forall \ell\in\mathcal{A}.
\end{equation}
Multiple states are constrained toward a common endpoint.

\paragraph{Adjacent relational self-distillation.}
\begin{equation}
    \mathbf{G}^{(\ell)}
    \rightarrow
    \mathbf{G}^{(\ell+1)}.
\end{equation}
Local student geometries are encouraged to remain consistent across depth.

\paragraph{GeoODE-KD.}
\begin{equation}
    \boxed{
    \mathbf{Z}^{(\ell+1)}
    \approx
    \operatorname{Exp}_{\mathbf{Z}^{(\ell)}}
    \left[
        \rho_\ell\,
        \mathbf{V}(
            \mathbf{Z}^{(\ell)},
            \mathbf{T}
        )
    \right]
    }.
\end{equation}

Here, the teacher specifies a \emph{directional transformation} at each
depth. The student is not required to arrive at the final teacher
representation prematurely; it is required to cover, at each depth, the
prescribed fraction of the remaining geodesic to the teacher, along a
direction that continuously decreases the teacher-conditioned potential.


% ============================================================
% OPTIONAL: END-OF-METHOD HYPOTHESIS PARAGRAPH
% ============================================================

\paragraph{Research hypothesis.}

Our central hypothesis is that a compact student benefits more from
learning a \emph{semantic trajectory} than from repeatedly imitating a
high-capacity teacher state. If correct, GeoODE-KD should exhibit three
observable behaviors. First, student-teacher semantic discrepancy should
decrease more smoothly through network depth than under static multi-layer
distillation; Corollary~1 makes this quantitative, predicting a geodesic
distance profile proportional to $R(t_\ell)$ rather than one that stays
flat and collapses in the last layers. Second, the teacher-student relational geometry gap should
contract progressively across layers rather than only at the output.
Third, these improvements should translate into better downstream
embedding quality without requiring stronger inference-time computation.
The experiments will directly test each of these hypotheses rather than
evaluating downstream accuracy alone.

\bibliographystyle{plainnat}
\bibliography{references}

\end{document}
